\chapter{Conclusion}
\label{chap:conclusion}
The overall project was a success. An affordable exoskeleton capable of virtual touch of physical objects was designed. The majority of the exoskeleton parts can be 3D printed, contributing to the ease of manufacturing and affordable price. The exoskeleton, besides providing force feedback, provides real-time information for the position and orientation of the exoskeleton. These exoskeleton features enable novel and intuitive ways to interact with the virtual environment. The exoskeleton has successfully aided in the training of the robot arm in the task of performing pick and place of soft objects in the simulation. It has done so by enabling expert trajectory collection in virtual environments. Expert trajectories are useful for training models that learn from human demonstrations, such as inverse reinforcement learning. While the training outcomes were only partially successful, the exoskeleton has demonstrated its potential. Beyond robot training, the exoskeleton's ability to convey feedback from virtual objects could benefit applications such as VR gaming, surgical simulations, and other research or virtual training applications. A future design could swap the sides of the brass inserts, making them easier to insert, resulting in a tighter design. The efficiency of the designed exoskeleton could be improved by manufacturing the double rack and pinion from metal instead of 3D printing it. Lastly, making the system wireless could improve the user experience and enable longer-distance movement.